% Сжатый текст статьи (Версия 2)
\documentclass[10pt,twocolumn,fleqn]{article}
\usepackage[utf8]{inputenc}
\usepackage[english,russian]{babel}
\usepackage{amsmath,amssymb,mathtools,geometry}
\geometry{a4paper,margin=1.2cm}
\DeclareMathOperator{\Rank}{Rank}

\title{\vspace{-4ex}\small\bf Автомат $7n$: Спектральные инварианты и тетрарные квазилогики}
\author{\small Группа $7n$}\date{\vspace{-9ex}}

\begin{document}
\maketitle

\begin{abstract}
\tiny Исследуется модель автомата $7n$ в терминах HoTT, алгебр Ли и проективных геометрий. На основе аксиомы относительности систем анализируются зазоры проективных плоскостей $N(6), N(10), N(12)$ по модулю 13.
\end{abstract}

\section{Онтология и аксиома относительности}
Число есть пересечение свойств контекста системы счета и относительных единиц ($o.e.$). \textit{Аксиома относительности:} любой элемент исследуется как единица объемлющей системы, либо как автономная система по отношению к своим подсистемам. Смена масштаба эквивалентна переходу h-уровней HoTT Воеводского. Булева логика $\mathbb{B}=\{0,1\}$ реализуется на полюсах фазовой базы $B=[0, \pi/2]$: $F_0 \simeq \mathbf{1}$ (День) и $F_{\pi/2} \simeq \mathbf{0}$ (Ночь). На экваторе $\theta = \pi/4$ деформация связности индуцирует тетрарную квазилогику.

\section{Микро- и мезомасштабные декомпозиции}
Числовые ритмы палеолита (Б.~Фролов) и логический календарь А.~Гротендника на 360 дней задают ограничения конечных групп. Порядок знакопеременной группы $A_6$ ($\vert A_6 \vert = 360$) декомпозируется через полиномы плоскостей $N(q) = q^2 + q + 1$:
\begin{equation}\small
    \vert A_6 \vert = N(12) + N(10) + N(6) + N(3) + \frac{3}{2}\vert S_4 \vert = 157 + 111 + 43 + 13 + 36
\end{equation}
Порядки $q=6,10$ запрещены (Брук--Райзер--Чоул), $q=12$ --- гипотетический открытый базис размерности $N(12)=157$. Топологический сдвиг $N(14) - N(12) = 54$ балансирует субфакториальный хаос: $!7 = 1854 = 1260 + (11 \times 54) = (42 \times 30) + (11 \times 9 \times \vert S_3 \vert)$. На отрезке $0 \dots 8$ произведение ячеек $\vert S_3 \vert \times \vert S_4 \vert = 144$ фиксирует константу чистых симметрий.

\section{Макромасштабный перенос Ли и спектр}
Редукция порядков $6\text{--}10\text{--}12 \pmod{13}$ ($N(3)=13$) задает октонионный перенос $\mathbb{O}$ оператором Ли $\nabla_{7n}(s, \theta)$ при инвариантности $\mathfrak{g}_2 \subset \mathfrak{so}(7)$ на тензоре $\text{Fano}_7$. Критическая точка фазового угла $\theta_{\text{crit}} = \arccos(156/157)$ для оператора $\hat{\mathbf{P}}^{(157)}$ формирует сингулярность: $\det(\hat{\mathbf{P}}^{(157)}(\theta_{\text{crit}})) = 0$. 

Кросс-масштабная ренормализация задана произведением $\hat{\mathbf{P}}_{\otimes} = \hat{\mathbf{P}}^{(7)} \otimes \hat{\mathbf{P}}^{(13)}$. При редукции граничного слоя по модулю 2 для $n=211$, Лапласиан равен $\hat{\mathbf{P}}_{\text{reg}}^{(211)} \equiv \hat{\mathbf{J}} - \hat{\mathbf{I}} \pmod 2$. Стабилизация ядра удерживается инвариантом максимума ранга $\Rank_{\mathbb{F}_2}(\hat{\mathbf{P}}_{\text{reg}}^{(211)}) = 210$, изолирующим расходимости хаоса субфакториала $!14 = 32\,071\,101\,049$.

\end{document}
